\documentclass{article}
\usepackage{graphicx} % Required for inserting images
\usepackage{amsmath}
\usepackage{amssymb}
\usepackage{bm}
\usepackage{fontawesome}
\usepackage{hyperref}
\usepackage{enumitem}
\usepackage{calc}
\usepackage{doi}
\usepackage[sorting=none]{biblatex}

\title{Comparative Ablation of Different SSM Architectures in DreamerV3}
\author{Dimitriye Danilović \and Shirisha Gujja \and Rohan Ohlan \and Nikil Thalapaneni}
\date{16 March 2026}
\addbibresource{refs.bib}
\begin{document}

\maketitle

\begin{abstract}
Since its introduction, DreamerV3\cite{dreamer} remains the foundational architecture for task-agnostic model-based reinforcement learning. Its RSSM component however is based on a GRU, which has known problems with long-range credit assignment and sequential training. In the time since, there has been much research into the use of different sequence model backbones, such as transformers or SSSMs. Most such attempts have not performed broad comparison across different task domains, with substantial variance in implementation other than the backbone. We present a comparative analysis of various backbones on identical tasks under task-agnostic configurations mirroring the evaluation methodology in the original paper by Hafner et al.
\end{abstract}

\section{Introduction}
Reinforcement learning has emerged as a powerful framework for training agents to perform complex tasks through interaction with their environment. 
\begin{itemize}
    \item \textbf{Model-free RL:} Methods that learn policies directly from experience without constructing an explicit model of the environment.
    \item \textbf{Model-based RL:} Methods that incorporate a learned model of the environment dynamics to plan and make decisions.
\end{itemize}

Prior to DreamerV3, the standard algorithm\cite{dreamer} for task-agnostic use was PPO\cite{ppo}, a model-free algorithm. However, PPO's applicability in task-agnostic contexts was limited, with many tasks demonstrating far higher performance from purpose-tuned expert models\cite{dreamer}. DreamerV3 introduced a model-based architecture with substantial gains in across-the-board performance with fixed hyperparameters, far exceeding that of PPO and on many tasks that of purpose-tuned expert models.

However, DreamerV3 continued to struggle with long-term extrapolation due to its architectural core\cite{memory}, a gated recurrent unit (GRU) model. GRUs have known difficulties with long-range credit assignment, as well as computational cost due to their sequential training. Since then, attempts have been made to replace the GRU as an architectural backbone in DreamerV3 with other types of sequence models:
\begin{enumerate}
    \item Samsami et al.\cite{memory} introduce Recall2Imagine (R2I) which utilizes a Structured State-Space Model in place of the GRU-based RSSM;
    \item Wang et al.\cite{drama} introduce Drama, which utilizes a Mamba-2\cite{mamba2} SSSM for efficient training;
    \item Mattes et al.\cite{hieros} introduce Hieros, which utilizes a Simplified Structured State-Space Model (S5)\cite{s5};
    \item Zhang et al.\cite{storm} introduce STORM, swapping in a transformer for the GRU;
    \item Burchi and Timofte\cite{twister} build on STORM with TWISTER, modifying the loss function to better accomodate the transformer swap.
\end{enumerate}

However, each of these works introduces additional architectural or methodological changes alongside the backbone swap, and evaluates on different subsets of environments. No prior work provides a controlled comparison isolating the sequence model as the sole variable. Herein we present such a controlled comparison.

\section{Related Work}
\subsection{Recall2Imagine --- S4-variant backbone}
Recall2Imagine\cite{memory} replaces the GRU with a multi-layer S4 variant (S3M), with parallel scan during training and sequential inference during imagination. This assists with long-range gradient flow and computational efficiency during training. However, they also make a change to the policy input, selecting from three different policy conditioning schemes based on environment type: output state policy \(\pi(a \mid z_t, h_t)\), hidden state policy \(\pi(a \mid z_t, x_t)\), or full state policy \(\pi(a \mid z_t, h_t, x_t)\). This per-domain tuning makes comparison ambiguous as it's unclear how much of the gain comes from the backbone vs. the policy conditioning choice. Per their ablation, \(\pi(a \mid z_t, h_t)\) under-performs on memory tasks, due to state compression. This is an architectural quality of SSSMs. Hence, for architectural uniformity, we elect to use \(\pi(a \mid z_t,x_t)\) for all SSSM-based models studied.
\subsection{Drama --- Mamba-2 variant backbone}
Drama\cite{drama} replaces the GRU with a Mamba-2 SSSM\cite{mamba2}, however, it also substantially restructures the overall model. DreamerV3's core novelty is the prior/posterior categorical latent with KL balancing, which Drama drops. The autoencoder architecture is also substantially altered, using a VAE trained on only the current observation, with no deterministic state input. Reward and termination predictors are single fully-connected layers from \(h_t\), rather than the MLP heads from concatenated \(z_t, h_t\) used in DreamerV3. They also introduce a new sampling methodology, dynamic frequency-based sampling, an orthogonal contribution which their ablation credits with nearly a quarter of the improvement. They also omit DreamerV3's robustness techniques for reward prediction.
\subsection{Hieros --- S5 backbone}
Hieros\cite{hieros} utilizes a S5\cite{s5} backend, but their core contribution is a goal-conditioned hierarchical policy structure. They ablate the hierarchical policy structure against the S5-backbone without it, but they do not broadly compare the S5-backbone alone with DreamerV3. They encountered similar issues with policy as R2I did.
\subsection{STORM --- Transformer backbone}
STORM\cite{storm} utilizes a GPT-style transformer\cite{gpt} backend. It's generally the closest in architecture of all backbone-swaps reviewed, with no significant differences other than the backbone.
\subsection{TWISTER --- Transformer backbone, optimized loss function}
TWISTER\cite{twister} utilizes the same transformer backend as STORM, but introduces a modified loss objective more well-suited for transformers. This is architecturally motivated, and not a typical confound, but useful to ablate in the overall comparison. It does utilize segment-level caching, which STORM does not, where hidden state sequences from previous segments are cached and reused as extended context. Notably however, they utilize relative positional encodings as introduced in Transformer-XL\cite{transformerxl} instead of absolute positional encodings as in STORM.

\section{Proposed Approach}
To harmonize the ablation, we will implement a unified architecture across backbones.

\subsection{SSSM Policy Structure}
All three SSSM approaches will be implemented using hidden state policy \(\pi(a \mid z_t, x_t)\) on all tasks. The standard DreamerV3 \(\pi(a \mid z_t, h_t)\) policy was demonstrated to under-perform significantly on memory tasks due to state compression which is inherent to the SSSM architecture by Samsami et al.\cite{memory} Other policy choices showed some gain on other tasks compared to the selected policy, however given the stated motivation of Dreamer-type models as task-agnostic with no change to hyperparameters, the choice of a single unified policy structure best suits the purposes of this analysis.

\subsection{Drop-in Mamba}
Notably, a drop-in replacement of the GRU with a Mamba-2 SSSM\cite{mamba2} will be performed, without the architectural differences posed by Drama\cite{drama}.

\subsection{Transformer Positional Encoding}
All transformer runs will be done with Transformer-XL-style\cite{transformerxl} relative positional encoding and segment-caching. The key difference between them will be the use of CPC or not.

\section{Experiment}
The core experiment will be the comparative ablation between the backbones with all else held equal, as well as the GRU DreamerV3 baseline. However, we also pose two sub-experiments:

\subsection{Dynamic Frequency-based Sampling}
The dynamic frequency-based sampling strategy used by Drama\cite{drama} demonstrated substantial improvement over uniform sampling of training data, which we implement for each model as a subexperiment. They maintain two per-transition count vectors:
\begin{itemize}
    \item \(v\) --- how many times a transition has been sampled for world model training,
    \item \(b\) --- how many times a transition has been sampled for imagination / policy training.
\end{itemize}

For world model sampling, the sampling probabilities are a simple softmax \(\sigma(-\mathbf{v})\). This is based on the anti-frequency sampling approach introduced by Robine et al.\cite{twm,drama}

For imagination sampling, the probabilities are given by
\[
\sigma(f(v, b)), \text{where } f(v, b) = v - b - \max(0, v - b).
\]

Thus, when \(v_i \ge b_i\) (the world model has seen this transition at least as often as the policy model has), \(f = 0\). When \(v_i < b_i\) (the policy model has seen this transition more than the world model), \(f = v_i - b_i\) which is negative, down-weighting the probability.

This serves to prevent training on transitions which the world-model has not yet learned to predict well.

\subsection{CPC}
TWISTER introduces a contrastive predictive coding loss motivated by the insight that next-state prediction does not fully exploit the capability of transformers\cite{twister}. They demonstrate a substantial gain over STORM with this methodology. Here, motivated by the Dao and Gu demonstration of the duality between transformers and SSSMs\cite{mamba2}, we implement such CPC loss to all models for comparison against baseline. We believe there may be substantial benefit to Mamba-2 and expect limited benefit to S4/S5.

The loss is calculated as follows:
\[
\mathcal{L}_{\mathrm{CPC}} = -\frac{1}{K}\sum_{k=0}^{K-1}\log\frac{\exp(\operatorname{sim}(z'_{t+k}, s_t))}{\sum_{z'_j\in Z'}\exp(\operatorname{sim}(z'_j, s_t))}
\]

\(z'\) denotes the encoding of augmented observation \(o'\) which is derived from \(o\) by randomly cropping and resizing.
The similarity function utilized is
\[
\operatorname{sim}(z'_j, s_t) = q^k_\phi(z'_j)^{T}p^k_\phi(s_t, a_{t:t+k})
\]
where \(q^k_\phi\) and \(p^k_\phi\) are two learned MLPs per each step k.

\section{Expected Deliverables}

\subsection{Code Repository}
A publicly available code repository implementing all backbone variants
(GRU, S4/S3M, S5, Mamba-2, Transformer, Transformer+CPC) within a unified
test harness built on the DreamerV3 codebase. The repository will include:
\begin{itemize}
    \item A shared training and evaluation pipeline with a single
          configuration interface for selecting backbone, environment
          suite, and sub-experiment options (dynamic frequency-based
          sampling, CPC loss).
    \item Reproduction scripts and configuration files sufficient to
          replicate all reported results.
    \item A verified GRU baseline reproducing published DreamerV3
          performance within expected variance on all three evaluation
          suites, serving as the reference point for all comparisons.
\end{itemize}

\subsection{Experimental Results}
Results from a minimum of 5 seeds per model--environment configuration,
reported as:
\begin{itemize}
    \item Mean and standard deviation of episodic return across seeds
          for each backbone on each evaluation suite (Atari100k, BSuite,
          POPGym).
    \item Interquartile mean (IQM) and performance profiles following
          the methodology recommended by Agarwal et al.\cite{rliable},
          to ensure statistically robust comparison across tasks and seeds.
    \item Per-environment learning curves at fixed evaluation intervals
          to characterize convergence behavior, not only final performance.
    \item Wall-clock training time per environment step for each backbone,
          measured on identical hardware, to quantify the practical
          efficiency implications of parallelizable (SSM, Transformer) vs.
          sequential (GRU) training.
\end{itemize}

\subsection{Ablation Analysis}
Controlled ablation results isolating the contribution of each
sub-experiment:
\begin{itemize}
    \item Dynamic frequency-based sampling vs.\ uniform sampling for
          each backbone, to determine whether the benefit demonstrated
          by Drama\cite{drama} generalizes across architectures.
    \item CPC loss vs.\ standard loss for each backbone, with particular
          attention to whether Mamba-2 benefits from CPC to a degree
          comparable to the transformer, as hypothesized from the
          transformer--SSM duality\cite{mamba2}. We will report the
          magnitude of CPC's effect on each backbone individually to
          allow fine-grained comparison.
\end{itemize}

\subsection{Paper}
A detailed paper documenting methodology, experimental results, and
analysis, including:
\begin{itemize}
    \item A discussion of any implementation decisions required to
          reconcile architectural differences between backbones within
          the unified harness, and their potential impact on results.
    \item Failure mode analysis for any backbone that substantially
          underperforms, examining whether the cause is architectural
          (e.g., state compression in SSSMs) or incidental (e.g.,
          hyperparameter sensitivity).
\end{itemize}

\subsection{Presentation}
A presentation summarizing findings, challenges encountered during
implementation, and proposed future research directions informed by
the results.

\section{Dataset and Evaluation Metrics}

\subsection{Evaluation Suites}
We evaluate across three environment suites, selected to test both
general control performance and memory-specific capabilities. This
mirrors the evaluation strategy of Samsami et al.\cite{memory}, who
demonstrated that memory-intensive benchmarks most clearly expose
the limitations of GRU-based world models.

\subsubsection{Atari100k}
The Atari100k benchmark\cite{atari,simple} restricts agents to 100k
environment interactions (400k frames with standard frame-skipping)
across the standard set of 26 games introduced by Kaiser et
al.\cite{simple}. This benchmark is used by all five motivating
backbone-swap papers, making it the natural common denominator for
comparison. It primarily tests sample-efficient visual control rather
than long-term memory, serving as our baseline for ensuring that
backbone changes do not regress general performance. We use sticky
actions ($p = 0.25$) following current best practices\cite{machado}.

\subsubsection{BSuite}
The Behaviour Suite for Reinforcement Learning\cite{bsuite} provides
targeted diagnostic experiments that isolate specific agent
capabilities. We focus on the memory-relevant subset:
\begin{itemize}
    \item \textbf{Memory Length}: a T-maze task parameterized by
          corridor length $N$, requiring the agent to recall a context
          bit presented at the start of the episode after $N$ steps of
          null observation. We evaluate across the full sweep of
          lengths to characterize the memory horizon of each backbone.
    \item \textbf{Memory Size}: extends Memory Length by requiring
          recall of multiple context bits simultaneously, testing
          memory capacity rather than just duration.
    \item \textbf{Discounting Chain}: a chain MDP of configurable
          length where the agent must propagate value over long
          horizons, directly testing credit assignment.
\end{itemize}
Following the protocol of Hafner et al.\cite{dreamer}, we run 10
seeds per BSuite configuration. BSuite environments use a single
environment instance by design.

\subsubsection{POPGym}
The Partially Observable Process Gym\cite{popgym} provides a diverse
collection of 15 POMDP environments with multiple difficulty levels.
Following the task selection of Samsami et al.\cite{memory} and Ni
et al., we evaluate on the three most memory-intensive environments:
\begin{itemize}
    \item \textbf{RepeatPrevious}: the agent must repeat an
          observation seen a fixed number of steps ago, directly
          testing sequential memorization.
    \item \textbf{Autoencode}: the agent must reproduce a sequence of
          observations after a delay, requiring multi-step memory
          encoding and retrieval.
    \item \textbf{Concentration}: a card-matching game requiring the
          agent to remember the locations and identities of
          previously revealed cards.
\end{itemize}
Each environment is evaluated at all three difficulty levels (Easy,
Medium, Hard), where difficulty scales the number of actions or
observations the agent must track simultaneously. Increased
difficulty requires correspondingly longer effective memory horizons,
providing a natural curriculum for distinguishing backbone capacity.

\subsection{Evaluation Metrics}
We adopt the evaluation methodology of Agarwal et al.\cite{rliable}
throughout, which has become the standard for statistically rigorous
comparison in deep RL.

\subsubsection{Score Normalization}
For Atari100k, raw scores are normalized to the human baseline
following the standard formula:
\[
    \text{HNS} = \frac{\text{agent score} -
    \text{random score}}{\text{human score} - \text{random score}}
\]
using the human and random baselines from Bellemare et
al.\cite{atari}. For BSuite, we use the native $[0,1]$ scoring
provided by the benchmark. For POPGym, we report raw episodic return,
as human baselines are not available for these environments.

\subsubsection{Aggregate Metrics}
We report the following aggregate statistics across tasks within each
suite:
\begin{itemize}
    \item \textbf{Interquartile Mean (IQM)}: the mean of scores
          between the 25th and 75th percentiles across tasks and
          seeds. IQM is more robust to outlier tasks than the mean
          while being more statistically efficient than the
          median\cite{rliable}.
    \item \textbf{Optimality Gap}: the amount by which the agent
          falls short of a reference threshold (1.0 HNS for
          Atari100k), averaged across tasks. This highlights
          underperformance on difficult tasks that IQM may obscure.
    \item \textbf{Performance Profiles}: cumulative distribution
          functions of per-task scores across all runs, following
          Agarwal et al.\cite{rliable}. These visualize the full
          distribution of performance and are more informative than
          any single scalar metric.
\end{itemize}
All aggregate metrics are reported with 95\% stratified bootstrap
confidence intervals computed over both tasks and seeds.

\subsubsection{Per-Environment Reporting}
In addition to aggregates, we report per-environment learning curves
at fixed evaluation intervals (every 10k steps for Atari100k, every
episode for BSuite). This serves two purposes: (1) characterizing
convergence speed, which may differ substantially between backbones
even when final performance is similar, and (2) identifying
environments where a particular backbone catastrophically fails,
which aggregate metrics can mask.

\subsubsection{Computational Efficiency}
We report wall-clock training time per environment step for each
backbone, measured on identical hardware (single H100 GPU). This
quantifies the practical benefit of parallelizable training
(SSMs, Transformers) over the sequential GRU, which is one of the
primary stated motivations for backbone replacement in the
literature\cite{memory,drama}.

\subsubsection{Statistical Testing}
For pairwise backbone comparisons, we compute the probability of
improvement\cite{rliable} --- the probability that a randomly
selected run-task pair from backbone $A$ outperforms one from
backbone $B$ --- as a nonparametric alternative to significance
testing. We consider a difference meaningful when the probability of
improvement exceeds 0.5 with non-overlapping 95\% confidence
intervals.

\section{Timeline}
To ensure consistent progress and a viable completion, the project will be carried out over a period of nine weeks and divided into the following phases:

\subsection{Weeks 1--2: Foundation \& Baselines:} Setting up the environment, getting acquainted with the DreamerV3 source, and reproducing the baseline GRU-based RSSM results on a chosen subset of environments.
\subsection{Weeks 3: Architecture Integration:} Putting the single code base into practice and combining the standardized hidden state policy with the S3M, Mamba-2, S5, and Transformer backbones.
\subsection{Weeks 4--5: Core Ablation Training:} Training and fine-tuning the implemented backbones in a sequential manner, monitoring the learning curves and confirming that the sequence models are operating appropriately within the task-agnostic constraints.
\subsection{Weeks 6--7: Sub-Experiments:} Implementation of Dynamic Frequency-based Sampling and Contrastive Predictive Coding (CPC) loss. Starting training runs for all architectures in these sub-experiments.
\subsection{Weeks 8: Evaluation \& Analysis:} Aggregating results, generating learning curves, and examining the impact of each backbone on long-range credit assignment and overall sample efficiency.
\subsection{Weeks 9: Deliverable Generation:} Finalizing the repository, drafting the final project report, and preparing the slide deck for the class presentation.


% \bibliographystyle{unsrt}
% \bibliography{refs}
\printbibliography


\end{document}